\section{Technical Overview}
\label{sec:technical-overview}

\textsc{Terrae} proves GBDT training by exposing the algebra already present in
histogram-based tree learning, rather than compiling the trainer into a large
generic circuit. 
In this section, we give a high-level overview of the main technical ingredients,
including how to batch claims across heterogeneous domains and how to prove histogram aggregation without a RAM argument.


\subsection{Batching Across Heterogeneous Domains}

Direct algebraic certification creates many relations over many different
domains. Training uses tensors indexed by rounds, classes, samples, nodes,
features, and bins; lookup, range, prefix-sum, division, maximum, and
zero-check relations therefore have different lengths and shapes. Ordinary
same-domain batching would force extra commitments or leave many small claims
unbatched. \textsc{Terrae} introduces two batching tools for this setting.

\paragraph{Domain-lifting batching.}
For a linear constraint over $\mathbb{H}_n$, after moving all terms to one side
we obtain a residual polynomial $F(X)$ and prove a divisibility claim
\[
    F(X)=Q(X)\nu_n(X),
    \qquad
    \nu_n(X)=X^n-1.
\]
Suppose two claims live over nested power-of-two domains
$\mathbb{H}_{n_1}\subset \mathbb{H}_{n_2}$. Instead of opening them separately,
\textsc{Terrae} lifts the smaller residual to the larger domain:
\[
    \widehat{F}_1(X)
    :=
    F_1(X)\frac{\nu_{n_2}(X)}{\nu_{n_1}(X)}.
\]
The verifier can then batch the lifted claim with the larger one as
\[
    \widehat{F}_1(X)+\eta F_2(X)=Q(X)\nu_{n_2}(X)
\]
for a random challenge $\eta$. This is used for zero-check and sum-check style
relations whose natural domains differ.

\paragraph{Interleaving batching.}
For local nonlinear checks, validity at one position depends only on the
witness values at that position. Examples include lookup, range, and
fixed-point division checks. For two vectors
$\mathbf{f},\mathbf{g}\in\mathbb{F}^n$ satisfying the same local relation,
\textsc{Terrae} forms the virtual interleaving
\[
    (\mathbf{f}[0],\mathbf{g}[0],\mathbf{f}[1],\mathbf{g}[1],\ldots).
\]
In polynomial form, if $f$ and $g$ are encoded over $\mathbb{H}_n$, the merged
object over $\mathbb{H}_{2n}$ is
\[
    \operatorname{Int}_n(f,g)(X)
    :=
    \operatorname{even}_n(X)f(X)
    +\operatorname{odd}_n(X)g(\omega_{2n}^{-1}X),
\]
where
\[
    \operatorname{even}_n(X)=\frac{X^n+1}{2},
    \qquad
    \operatorname{odd}_n(X)=\frac{1-X^n}{2}.
\]
For different vector lengths, the shorter vector is first aligned by a domain
substitution $X\mapsto X^r$, which replicates it over the larger domain. The
merged vectors are virtual: their openings are derived from existing
commitments, public sparse selectors, and domain substitutions, so batching
does not require committing to new auxiliary witness polynomials.

Together, these ingredients are the main contribution of \textsc{Terrae}: the
training trace remains explicit enough to express the semantics of GBDT
training, while the proof system avoids both a monolithic training circuit and
a large collection of separately opened auxiliary commitments. The next section
describes how the committed forest, inference proof, well-formedness checks,
training proof, and split-selection relations are assembled.


\subsection{Histogram Aggregation Without RAM}

GBDT training repeatedly converts sample-wise information into per-node,
per-feature, and per-bin statistics. Prior certification-style systems~\cite{sparrow24,zkboost26}
naturally model this as a memory update process: for each routed sample, update
the corresponding histogram bucket and prove the resulting RAM trace.
\textsc{Terrae} instead uses a histogram-specific argument.

For simplicity, let's consider a simplified version of histogram aggregation.
Suppose we have vectors $\mathbf{ind}\in [m]^n$ and $\mathbf{u}, \mathbf{v}\in \mathbb{F}^n$.
Our goal is to prove that $\mathbf{v}$ is the histogram
of $\mathbf{u}$ according to the indices in $\mathbf{ind}$:
\begin{align}
    \mathbf{v}[j]=\sum_{i:\mathbf{ind}[i]=j}\mathbf{u}[i].
\end{align}

To prove this, the verifier generates a random challenge
$\alpha\gets \mathbb{F}$, which binds each $i$-th slot to $\alpha^i$. The parties next move to prove
\begin{align}
    \sum_{i=0}^{n-1}\mathbf{u}[i]\cdot \alpha^{\mathbf{ind}[i]}
    =
    \sum_{j=0}^{m-1}\mathbf{v}[j]\cdot \alpha^j.\label{eq:histogram-aggregation-power-sum}
\end{align}
The remaining question is how to prove the exponentiated terms
$\alpha^{\mathbf{ind}[i]}$.
The prover commits a power table
$\mathbf{p}\in\mathbb{F}^{m}$ with intended entries
$\mathbf{p}[j]=\alpha^j$. It proves that $\mathbf{p}$ is well formed by
checking
\begin{align}
    \mathbf{p}[0]&=1,\\
    \mathbf{p}[j]&=\alpha\cdot \mathbf{p}[j-1],
    \qquad j\in\{1,\ldots,m-1\}.
\end{align}
Or equivalently in polynomial form,
\begin{align}
    &p(\omega_m^0) = 1,\\
    &(p(t) - \alpha\cdot p(\omega_m^{-1} t))\cdot (1 - \lambda_m^{(0)}(t)) = q(X)\nu_{m}(X),
\end{align}
where recall $\lambda_m^{(0)}(X)$ is a Lagrange
basis polynomial, and the second relation 
is a zero-check.

The prover then proves by lookup that each row
$(\mathbf{ind}[i],\mathbf{pow}[i])$ appears in the public-indexed table
$(j,\mathbf{pow}[j])_{j\in[m]}$, i.e.
\begin{align}
    \mathbf{ind}\| \mathbf{pow}\subseteq (0, 1, \dots, m-1)^\top\| \mathbf{p}.
\end{align}
This binds
$\mathbf{pow}[i]=\alpha^{\mathbf{ind}[i]}$, after which
\Cref{eq:histogram-aggregation-power-sum}
becomes
\begin{align}
    \sum_{i=0}^{n-1}\mathbf{u}[i]\cdot \mathbf{pow}[i]
    =
    \sum_{j=0}^{m-1}\mathbf{v}[j]\cdot \mathbf{p}[j],
\end{align}
which can be checked by standard sum-check.

Soundness follows from the randomness of $\alpha$. 
To analyze the soundness error introduced by
\Cref{eq:histogram-aggregation-power-sum},
define
\[
    \Delta_j
    =
    \mathbf{v}[j]
    -
    \sum_{i:\mathbf{ind}[i]=j}\mathbf{u}[i].
\]
If the claim that $\mathbf{v}$ is the histogram of $\mathbf{u}$ is false, 
then some $\Delta_j\neq 0$, and the accepted aggregate identity implies
\[
    \sum_{j=0}^{m-1}\Delta_j\alpha^j=0.
\]
This is a nonzero polynomial in $\alpha$ of degree at most $m-1$, so a cheating
prover can satisfy it with probability at most $(m-1)/|\mathbb{F}|$, assuming
$\alpha$ is sampled after the relevant commitments are fixed. Thus the
aggregate identity is a compact randomized check of the full histogram relation.